\documentclass[12pt,a4paper]{article}

% Codificação e idioma
\usepackage[utf8]{inputenc}
\usepackage[portuguese]{babel}

% Pacotes úteis
\usepackage{booktabs}
\usepackage{longtable}
\usepackage{graphicx}
\usepackage{hyperref}
\usepackage{geometry}
\usepackage{indentfirst}
\usepackage{float}

% Margens
\geometry{margin=2.5cm}


\newcommand{\projetotitulo}{Introdução à Manufatura Aditiva}
\newcommand{\grupoid}{Grupo Manufatura Aditiva}
\newcommand{\sprintnum}{4}
\newcommand{\sprintperiodo}{08/06 a 21/06/2026}
\newcommand{\scrummaster}{Lucas Moura}
\newcommand{\gitrepo}{https://github.com/viniciusnovelo/Introducao-a-Manufatura-Aditiva}
\newcommand{\gitcommit}{} % hash do commit desta entrega

\begin{document}

\begin{center}
{\Large\bfseries Relatório de Sprint \sprintnum}\\
{\large \projetotitulo}\\
\grupoid{} \ | \ \sprintperiodo{} \ | \ SM: \scrummaster
\end{center}

\section{Sprint Backlog}
\footnote{hash: \gitcommit}

\begin{longtable}{|c|p{3.8cm}|p{1.9cm}|p{2.2cm}|p{1.7cm}|p{4.2cm}|}
\toprule
ID & História & Estimativa & Responsável & Status & Critério de Aceitação \\
\midrule
1 & Criação do Instagram para Divulgação & Sprint 1 & Lucas Moura & Concluído & Perfil criado, configurado com identidade visual do projeto e contendo ao menos uma publicação inicial. \\ \hline

2 & Criação do formulário de divulgação & Sprint 1 & Vinícius & Concluído & Formulário criado, testado e apto para receber inscrições dos participantes. \\ \hline

3 & Início da divulgação para comunidade interna e externa & Sprint 1 & Vinícius & Concluído & Divulgação realizada em ao menos uma escola ou meio de comunicação definido pela equipe. \\ \hline

4 & Criação de slides para as aulas de SolidEdge para os calouros & Sprint 1 & Paulo & Concluído & Slides elaborados e revisados pela equipe antes da aplicação das aulas. \\ \hline

5 & Criação do quadro Kanban & Sprint 1 & Lucas Leão & Concluído & Quadro criado contendo colunas de acompanhamento e tarefas cadastradas. \\ \hline

6 & Coleta e organização de dados das inscrições dos alunos & Sprint 1 & Paulo & Concluído & Dados armazenados e organizados em planilha. \\ \hline

7 & Agendamento das salas e laboratórios que serão utilizados & Sprint 1 & Leonardo & Concluído & Reservas confirmadas junto à instituição para todas as aulas previstas. \\ \hline

8 & Verificação e manutenção das impressoras 3D e softwares no laboratório & Sprint 1 & Yuri & Concluído & Equipamentos testados e softwares funcionando corretamente antes do início do curso. \\ \hline

9 & Criação de relatório de Sprint 1 & Sprint 1 & Lucas Leão & Concluído & Relatório elaborado contendo atividades realizadas, resultados e pendências. \\ \hline

10 & Criação do repositório no Github & Sprint 1 & Vinícius & Concluído & Repositório criado com acesso para a equipe e estrutura inicial organizada. \\ \hline

11 & Treinamento da equipe nas ferramentas de modelagem 3D & Sprint 1 e 2 & Lucas Moura & Concluído & Equipe capacitada para executar operações básicas de modelagem no software. \\ \hline

12 & Elaboração dos materiais para as aulas 1 e 2 do curso & Sprint 2 & Equipe & Concluído & Materiais preparados e revisados antes da realização das aulas. \\ \hline

13 & Realização da aula 1 & Sprint 2 & Equipe & Concluído & Aula ministrada conforme cronograma e com registro de presença dos alunos. \\ \hline

14 & Realização da aula 2 & Sprint 2 & Equipe & Concluído & Aula ministrada conforme cronograma e com registro de presença dos alunos. \\ \hline

15 & Criação de relatório de Sprint 2 & Sprint 2 & Leonardo, Paulo, Yuri & Concluído & Relatório elaborado contendo atividades realizadas, resultados e pendências. \\ \hline

16 & Atualização do quadro Kanban & Sprint 2 & Lucas Leão & Concluído & Todas as tarefas da sprint atualizadas conforme o andamento das atividades. \\ \hline

17 & Atualização do repositório Github & Sprint 2 & Vinícius & Concluído & Materiais e documentos da sprint enviados e organizados no repositório. \\ \hline

18 & Treinamento da equipe nas ferramentas de impressão 3D & Sprint 3 & Lucas Moura & Concluído & Equipe apta a operar as impressoras e configurar impressões básicas. \\ \hline

19 & Elaboração dos materiais para as aulas 3 e 4 do curso & Sprint 3 & Equipe & Concluído & Materiais preparados e aprovados pela equipe antes das aulas. \\ \hline

20 & Realização da aula 3 & Sprint 3 & Equipe & Concluído & Aula ministrada conforme cronograma e com registro de presença dos alunos. \\ \hline

21 & Criação de relatório de Sprint 3 & Sprint 3 & Leonardo, Paulo, Yuri & Concluído & Relatório elaborado contendo atividades realizadas, resultados e pendências. \\ \hline

22 & Atualização do quadro Kanban & Sprint 3 & Lucas Leão & Concluído & Todas as tarefas da sprint atualizadas conforme o andamento das atividades. \\ \hline

23 & Atualização do repositório Github & Sprint 3 & Vinícius & Concluído & Materiais e documentos da sprint enviados e organizados no repositório. \\ \hline

24 & Elaboração dos materiais para as aulas 5 e 6 do curso & Sprint 4 & Lucas Moura & Concluído & Materiais finalizados, revisados e disponíveis para aplicação das aulas. \\ \hline

25 & Realização da aula 4 & Sprint 4 & Equipe & Concluído & Aula ministrada conforme cronograma e com registro de presença dos alunos. \\ \hline

26 & Realização da aula 5 & Sprint 4 & Equipe & Concluído & Aula ministrada conforme cronograma e com registro de presença dos alunos. \\ \hline

27 & Criação de relatório de Sprint 4 & Sprint 4 & Leonardo, Paulo, Yuri & Concluído & Relatório elaborado contendo atividades realizadas, resultados e pendências. \\ \hline

28 & Atualização do quadro Kanban & Sprint 4 & Lucas Leão & Em andamento & Todas as tarefas da sprint atualizadas conforme o andamento das atividades. \\ \hline

29 & Atualização do repositório Github & Sprint 4 & Vinícius & Em andamento & Materiais e documentos da sprint enviados e organizados no repositório. \\ \hline

30 & Realização da aula 6 & Sprint 5 & Equipe & Não iniciado & Aula final realizada e conteúdo previsto concluído. \\ \hline

31 & Criação de formulário para coleta de feedback & Sprint 5 & Leonardo, Yuri & Não iniciado & Formulário criado contendo perguntas objetivas e discursivas sobre o curso. \\ \hline

32 & Coleta de feedback dos alunos através de formulário & Sprint 5 & Leonardo, Yuri & Não iniciado & Respostas coletadas de parte significativa dos participantes do curso. \\ \hline

33 & Criação de relatório de Sprint 5 & Sprint 5 & Leonardo, Paulo, Yuri & Não iniciado & Relatório elaborado contendo atividades realizadas, resultados e pendências. \\ \hline

34 & Atualização do quadro Kanban & Sprint 5 & Lucas Leão & Não iniciado & Todas as tarefas da sprint atualizadas conforme o andamento das atividades. \\ \hline

35 & Atualização do repositório Github & Sprint 5 & Vinícius & Não iniciado & Materiais e documentos da sprint enviados e organizados no repositório. \\ \hline

36 & Criação do relatório final & Sprint 5 & A definir & Não iniciado & Relatório final concluído contendo descrição completa do projeto, resultados e conclusões. \\

\bottomrule
\end{longtable}

\section{Burndown Chart}
\begin{figure}[h!]
\centering
\includegraphics[width=1.0\textwidth]{Imagens/burndown_sprint\sprintnum.png}}
\caption{Burndown -- Sprint \sprintnum}
\end{figure}

% Análise do burndown (mínimo 3 linhas):
O gráfico acima representa o engajamento e o trabalho efetivo realizado pelo grupo durante a Sprint 4, sendo que a linha azul representa o progresso ideal no decorrer de cada dia e a linha laranja o progresso realizado efetivamente. Para cada atividade da Sprint 4 foi atribuído um valor de esforço representado no eixo vertical do gráfico.

\section{Artefatos Produzidos}

Em continuidade aos materiais desenvolvidos nas semanas anteriores, foi finalizado os materiais para as aulas do curso, com a conclusão dos materiais referentes às aulas 5 e 6. O que inclui apresentações em Power Point, design de peças no Solidge Edge para prática e apostilas para introdução na impressão 3D.
\begin{figure}[H]
\centering

\begin{minipage}{0.48\textwidth}
    \centering
    \includegraphics[width=\textwidth]{Imagens/Slide Aula 5.png}
    \caption{Slide - Aula 5}
\end{minipage}
\hfill
\begin{minipage}{0.48\textwidth}
    \centering
    \includegraphics[width=\textwidth]{Imagens/Slide Impressão 3D.png}
    \caption{Slide - Guia Impressão 3D}
\end{minipage}

\vspace{2.0cm}

\begin{minipage}{0.58\textwidth}
    \centering
    \includegraphics[width=\textwidth]{Imagens/Apostila.png}
        \caption{Apostila - Configurações Impressora 3D}
\end{minipage}
\end{figure}


\section{Atualização com o Parceiro}
O contato com os parceiros continuou de forma contínua via WhatsApp e encontros presenciais. No período, foram ministradas com sucesso as Aulas 4 e 5 do cronograma. A Aula 4 consistiu na continuação direta da montagem das peças no software pelos alunos, reaproveitando o material didático da Aula 3. Já a Aula 5 foi integralmente dedicada às atividades práticas de revolução e à introdução aos parâmetros de fatiamento e impressão 3D.

\section{Participação do Calouro}
Os calouros demonstraram excelente autonomia e engajamento prático no desenvolvimento de materiais de suporte e durante as aulas presenciais. Na Aula 5, eles atuaram diretamente no suporte aos alunos, auxiliando e orientando a execução correta dos exercícios práticos de modelagem geométrica e revolução aplicados no desenvolvimento das peças.

\section{Impedimentos e Desvios}
Os impedimentos críticos apontados na sprint anterior foram mitigados com sucesso pela equipe. O problema com a expiração das licenças do Solid Edge foi normalizado, o que viabilizou o avanço regular e a conclusão das etapas de montagem de peças programadas para a Aula 4. Não houve novos impedimentos ou atrasos nas duas últimas semanas.

\section{Retrospectiva}
\subsection*{O que funcionou bem}
A ação concreta estipulada na última sprint foi totalmente cumprida com a entrega do manual didático de fatiamento no XYZware Pro e da tabela com os parâmetros ideais de impressão para os materiais ABS, PETG e PLA. O guia instruiu o preenchimento dos separadores de Supports (habilitando suporte obrigatório) , General (altura de camada) , Infill (densidade em 20\% com padrão Colmeia/Honeycomb) e Shells. O rendimento dos alunos na modelagem por revolução na Aula 5 também atendeu às expectativas. 

\subsection*{O que precisa melhorar}
É necessário otimizar o tempo de acompanhamento da primeira camada da impressão em sala , visto que o aquecimento completo da mesa (plataforma aquecida) costuma demandar um tempo considerável, especialmente para o filamento ABS, cuja parametrização exige metas térmicas elevadas entre 90°C e 110°C. 

\subsection*{Ação concreta para a próxima sprint}
Garantir a aplicação correta dos métodos de adesão na mesa de alumínio para evitar que a fundação descole e a impressão falhe. Os alunos serão ensinados a cobrir a mesa com fitas adesivas (Bed Tapes, fita crepe azul larga ou fita Kapton) e aplicar uma fina camada de cola bastão ou spray fixador de cabelo de fixação extra forte antes do início de cada ciclo físico. 

\section{Planejamento da Próxima Sprint}
Na próxima sprint , a equipe dará prosseguimento às etapas práticas de manufatura aditiva utilizando o software XYZware Pro. O objetivo primordial será guiar os alunos no processo físico de calibração e nivelamento manual da mesa acessando o menu UTILITIES > CALIBRATE no painel de LCD. Os alunos aprenderão a mecânica dos "Passos" utilizando os três botões brancos (knobs) debaixo da mesa , sabendo que cada clique sentido equivale a 90 graus de rotação (1 Passo) , e corrigindo a mesa conforme as direções exibidas em tela (como Left Knob, Front Knob ou Right Knob) até obter a validação de "PERFECT". Também será fixada uma rotina de manutenção preventiva instruindo os alunos a limparem o sensor metálico ao lado do bico aquecido com a escova de aço antes de calibrar, pois resíduos de plástico ressecado fazem o contato elétrico falhar.
\end{document}